% Common packages for the Chinese experiment report.
% Compile with XeLaTeX.
\usepackage{fontspec}
\usepackage{xeCJK}
\usepackage{microtype}
\usepackage[letterpaper,margin=1in]{geometry}
\usepackage{graphicx}
\usepackage{amsmath,amssymb,mathbbol}
\usepackage{acronym}
\usepackage{enumitem}
\usepackage[pagebackref=true,breaklinks=true,colorlinks]{hyperref}
\usepackage{balance}
\usepackage{xspace}
\usepackage{setspace}
\usepackage{indentfirst}
\usepackage[skip=3pt,font=small]{subcaption}
\usepackage[skip=3pt,font=small]{caption}
\usepackage[dvipsnames]{xcolor}
\usepackage[capitalise]{cleveref}
\usepackage{booktabs,tabularx,colortbl,multirow,array,makecell}
\usepackage{siunitx}
\usepackage{float}

\IfFontExistsTF{Times New Roman}{\setmainfont{Times New Roman}}{\setmainfont{TeX Gyre Termes}}
\IfFontExistsTF{Helvetica}{\setsansfont{Helvetica}}{\setsansfont{TeX Gyre Heros}}
\IfFontExistsTF{Menlo}{\setmonofont{Menlo}}{\setmonofont{Inconsolata}}
\IfFontExistsTF{Songti SC}{
  \setCJKmainfont{Songti SC}[AutoFakeBold=2]
  \setCJKsansfont{Songti SC}[AutoFakeBold=2]
  \setCJKmonofont{Songti SC}[AutoFakeBold=2]
}{
  \setCJKmainfont{FandolSong-Regular.otf}[AutoFakeBold=2]
  \setCJKsansfont{FandolHei-Regular.otf}
  \setCJKmonofont{FandolSong-Regular.otf}
}

\hypersetup{pdfencoding=auto,colorlinks=true,allcolors=black}
\pagestyle{plain}

% Chinese document names.
\renewcommand{\figurename}{图}
\renewcommand{\tablename}{表}
\renewcommand{\refname}{参考资料}
\renewenvironment{abstract}%
{%
  \vskip 0.075in%
  \centerline%
  {\large\bf 摘要}%
  \vspace{0.5ex}%
  \begin{quote}%
  \setlength{\parindent}{2em}%
  \leavevmode\hspace*{2em}\ignorespaces%
}
{
  \par%
  \end{quote}%
  \vskip 1ex%
}

% Handy shorthand
\makeatletter
\DeclareRobustCommand\onedot{\futurelet\@let@token\@onedot}
\def\@onedot{\ifx\@let@token.\else.\null\fi\xspace}
\def\eg{\emph{e.g}\onedot}
\def\Eg{\emph{E.g}\onedot}
\def\ie{\emph{i.e}\onedot}
\def\Ie{\emph{I.e}\onedot}
\def\cf{\emph{c.f}\onedot}
\def\Cf{\emph{C.f}\onedot}
\def\etc{\emph{etc}\onedot}
\def\vs{\emph{vs}\onedot}
\def\wrt{w.r.t\onedot}
\def\dof{d.o.f\onedot}
\def\etal{\emph{et al}\onedot}
\makeatother

\definecolor{gray}{gray}{0.9}
\definecolor{reportgreen}{HTML}{008E7D}
\definecolor{reportorange}{HTML}{D97706}
\definecolor{reportslate}{HTML}{334155}

% Table style
\sisetup{
  detect-all,
  table-number-alignment = center,
  group-minimum-digits = 4
}
\newcolumntype{Y}{>{\raggedright\arraybackslash}X}
\newcolumntype{L}[1]{>{\raggedright\arraybackslash}p{#1}}
\renewcommand{\arraystretch}{1.12}

% Handy math ops
\DeclareMathOperator*{\argmax}{arg\,max}
\DeclareMathOperator*{\argmin}{arg\,min}
\newcommand{\norm}[1]{\left\Vert #1 \right\Vert}

% Spacing
\frenchspacing
\setstretch{1}
\setlength{\parindent}{2em}
\setlength{\parskip}{5.5pt}
\makeatletter
\renewcommand{\normalsize}{%
  \@setfontsize\normalsize\@xpt{13.75pt}%
  \abovedisplayskip      7\p@ \@plus 2\p@ \@minus 5\p@
  \abovedisplayshortskip \z@ \@plus 3\p@
  \belowdisplayskip      \abovedisplayskip
  \belowdisplayshortskip 4\p@ \@plus 3\p@ \@minus 3\p@
}
\normalsize
\makeatother
\setlength\floatsep{0.5\baselineskip plus 3pt minus 2pt}
\setlength\textfloatsep{0.5\baselineskip plus 3pt minus 2pt}
\setlength\dbltextfloatsep{0.5\baselineskip plus 3pt minus 2pt}
\setlength\intextsep{0.5\baselineskip plus 3pt minus 2pt}

\makeatletter
\renewcommand{\maketitle}{%
  \begin{center}
    {\LARGE\bfseries \@title\par}
    \vspace{1.2em}
    {\normalsize \@author\par}
  \end{center}
  \vspace{1.2em}
}
\makeatother

\makeatletter
\renewcommand{\paragraph}{%
  \@startsection{paragraph}{4}%
  {\z@}{0ex \@plus 0ex \@minus 0ex}{-1em}%
  {\hskip\parindent\normalfont\normalsize\bfseries}%
}
\makeatother

% Graphics path
\graphicspath{{figures/}}

% Clever references
\crefname{algorithm}{算法}{算法}
\Crefname{algorithm}{算法}{算法}
\crefname{section}{节}{节}
\Crefname{section}{节}{节}
\crefname{table}{表}{表}
\Crefname{table}{表}{表}
\crefname{figure}{图}{图}
\Crefname{figure}{图}{图}

% Acronym
\acrodef{ttft}[TTFT]{Time To First Token}
\acrodef{rss}[RSS]{Resident Set Size}
